\begin{register}{H}{UART\_FIFO\_REG}{0x0}\label{regdesc:UARTFIFOREG}
\regfield{(reserved)}{24}{8}{000000000000000000000000}%
\regfield{UART\_RXFIFO\_RD\_BYTE}{8}{0}{00000000}%
\reglabel{Reset}\regnewline%
\begin{regdesc}\begin{reglist}
\item [UART\_RXFIFO\_RD\_BYTE] 访问 FIFO。（读/写）
\end{reglist}\end{regdesc}
\end{register}


\begin{register}{H}{UART\_INT\_RAW\_REG}{0x4}\label{regdesc:UARTINTRAWREG}
\regfield{(reserved)}{13}{19}{0000000000000}%
\regfield{UART\_AT\_CMD\_CHAR\_DET\_INT\_RAW}{1}{18}{0}%
\regfield{UART\_RS485\_CLASH\_INT\_RAW}{1}{17}{0}%
\regfield{UART\_RS485\_FRM\_ERR\_INT\_RAW}{1}{16}{0}%
\regfield{UART\_RS485\_PARITY\_ERR\_INT\_RAW}{1}{15}{0}%
\regfield{UART\_TX\_DONE\_INT\_RAW}{1}{14}{0}%
\regfield{UART\_TX\_BRK\_IDLE\_DONE\_INT\_RAW}{1}{13}{0}%
\regfield{UART\_TX\_BRK\_DONE\_INT\_RAW}{1}{12}{0}%
\regfield{UART\_GLITCH\_DET\_INT\_RAW}{1}{11}{0}%
\regfield{UART\_SW\_XOFF\_INT\_RAW}{1}{10}{0}%
\regfield{UART\_SW\_XON\_INT\_RAW}{1}{9}{0}%
\regfield{UART\_RXFIFO\_TOUT\_INT\_RAW}{1}{8}{0}%
\regfield{UART\_BRK\_DET\_INT\_RAW}{1}{7}{0}%
\regfield{UART\_CTS\_CHG\_INT\_RAW}{1}{6}{0}%
\regfield{UART\_DSR\_CHG\_INT\_RAW}{1}{5}{0}%
\regfield{UART\_RXFIFO\_OVF\_INT\_RAW}{1}{4}{0}%
\regfield{UART\_FRM\_ERR\_INT\_RAW}{1}{3}{0}%
\regfield{UART\_PARITY\_ERR\_INT\_RAW}{1}{2}{0}%
\regfield{UART\_TXFIFO\_EMPTY\_INT\_RAW}{1}{1}{0}%
\regfield{UART\_RXFIFO\_FULL\_INT\_RAW}{1}{0}{0}%
\reglabel{Reset}\regnewline%
\begin{regdesc}\begin{reglist}
\item [UART\_AT\_CMD\_CHAR\_DET\_INT\_RAW] \hyperref[int:UARTATCMDCHARDETINT]{UART\_AT\_CMD\_CHAR\_DET\_INT} 中断的原始中断状态位。（只读）
\item [UART\_RS485\_CLASH\_INT\_RAW] \hyperref[int:UARTRS485CLASHINT]{UART\_RS485\_CLASH\_INT} 中断的原始中断状态位。（只读）
\item [UART\_RS485\_FRM\_ERR\_INT\_RAW] \hyperref[int:UARTRS485FRMERRINT]{UART\_RS485\_FRM\_ERR\_INT} 中断的原始中断状态位。（只读）
\item [UART\_RS485\_PARITY\_ERR\_INT\_RAW] \hyperref[int:UARTRS485PARITYERRINT]{UART\_RS485\_PARITY\_ERR\_INT} 中断的原始中断状态位。（只读）
\item [UART\_TX\_DONE\_INT\_RAW] \hyperref[int:UARTTXDONEINT]{UART\_TX\_DONE\_INT} 中断的原始中断状态位。（只读）
\item [UART\_TX\_BRK\_IDLE\_DONE\_INT\_RAW] \hyperref[int:UARTTXBRKIDLEDONEINT]{UART\_TX\_BRK\_IDLE\_DONE\_INT} 中断的原始中断状态位。（只读）
\item [UART\_TX\_BRK\_DONE\_INT\_RAW] \hyperref[int:UARTTXBRKDONEINT]{UART\_TX\_BRK\_DONE\_INT} 中断的原始中断状态位。（只读）
\item [UART\_GLITCH\_DET\_INT\_RAW] \hyperref[int:UARTGLITCHDETINT]{UART\_GLITCH\_DET\_INT} 中断的原始中断状态位。（只读）
\item [UART\_SW\_XOFF\_INT\_RAW] \hyperref[int:UARTSWXOFFINT]{UART\_SW\_XOFF\_INT} 中断的原始中断状态位。（只读）
\item [UART\_SW\_XON\_INT\_RAW] \hyperref[int:UARTSWXONINT]{UART\_SW\_XON\_INT} 中断的原始中断状态位。（只读）
\item [UART\_RXFIFO\_TOUT\_INT\_RAW] \hyperref[int:UARTRXFIFOTOUTINT]{UART\_RXFIFO\_TOUT\_INT} 中断的原始中断状态位。（只读）
\item [UART\_BRK\_DET\_INT\_RAW] \hyperref[int:UARTBRKDETINT]{UART\_BRK\_DET\_INT} 中断的原始中断状态位。（只读）
\item [UART\_CTS\_CHG\_INT\_RAW] \hyperref[int:UARTCTSCHGINT]{UART\_CTS\_CHG\_INT} 中断的原始中断状态位。（只读）
\item [UART\_DSR\_CHG\_INT\_RAW] \hyperref[int:UARTDSRCHGINT]{UART\_DSR\_CHG\_INT} 中断的原始中断状态位。（只读）
\item [UART\_RXFIFO\_OVF\_INT\_RAW] \hyperref[int:UARTRXFIFOOVFINT]{UART\_RXFIFO\_OVF\_INT} 中断的原始中断状态位。（只读）
\item [UART\_FRM\_ERR\_INT\_RAW] \hyperref[int:UARTFRMERRINT]{UART\_FRM\_ERR\_INT} 中断的原始中断状态位。（只读）
\item [UART\_PARITY\_ERR\_INT\_RAW] \hyperref[int:UARTPARITYERRINT]{UART\_PARITY\_ERR\_INT} 中断的原始中断状态位。（只读）
\item [UART\_TXFIFO\_EMPTY\_INT\_RAW] \hyperref[int:UARTTXFIFOEMPTYINT]{UART\_TXFIFO\_EMPTY\_INT} 中断的原始中断状态位。（只读）
\item [UART\_RXFIFO\_FULL\_INT\_RAW] \hyperref[int:UARTRXFIFOFULLINT]{UART\_RXFIFO\_FULL\_INT} 中断的原始中断状态位。（只读）
\end{reglist}\end{regdesc}
\end{register}


\begin{register}{H}{UART\_INT\_ST\_REG}{0x8}\label{regdesc:UARTINTSTREG}
\regfield{(reserved)}{13}{19}{0000000000000}%
\regfield{UART\_AT\_CMD\_CHAR\_DET\_INT\_ST}{1}{18}{0}%
\regfield{UART\_RS485\_CLASH\_INT\_ST}{1}{17}{0}%
\regfield{UART\_RS485\_FRM\_ERR\_INT\_ST}{1}{16}{0}%
\regfield{UART\_RS485\_PARITY\_ERR\_INT\_ST}{1}{15}{0}%
\regfield{UART\_TX\_DONE\_INT\_ST}{1}{14}{0}%
\regfield{UART\_TX\_BRK\_IDLE\_DONE\_INT\_ST}{1}{13}{0}%
\regfield{UART\_TX\_BRK\_DONE\_INT\_ST}{1}{12}{0}%
\regfield{UART\_GLITCH\_DET\_INT\_ST}{1}{11}{0}%
\regfield{UART\_SW\_XOFF\_INT\_ST}{1}{10}{0}%
\regfield{UART\_SW\_XON\_INT\_ST}{1}{9}{0}%
\regfield{UART\_RXFIFO\_TOUT\_INT\_ST}{1}{8}{0}%
\regfield{UART\_BRK\_DET\_INT\_ST}{1}{7}{0}%
\regfield{UART\_CTS\_CHG\_INT\_ST}{1}{6}{0}%
\regfield{UART\_DSR\_CHG\_INT\_ST}{1}{5}{0}%
\regfield{UART\_RXFIFO\_OVF\_INT\_ST}{1}{4}{0}%
\regfield{UART\_FRM\_ERR\_INT\_ST}{1}{3}{0}%
\regfield{UART\_PARITY\_ERR\_INT\_ST}{1}{2}{0}%
\regfield{UART\_TXFIFO\_EMPTY\_INT\_ST}{1}{1}{0}%
\regfield{UART\_RXFIFO\_FULL\_INT\_ST}{1}{0}{0}%
\reglabel{Reset}\regnewline%
\begin{regdesc}\begin{reglist}
\item [UART\_AT\_CMD\_CHAR\_DET\_INT\_ST] \hyperref[int:UARTATCMDCHARDETINT]{UART\_AT\_CMD\_CHAR\_DET\_INT} 中断的隐蔽中断状态位。（只读）
\item [UART\_RS485\_CLASH\_INT\_ST] \hyperref[int:UARTRS485CLASHINT]{UART\_RS485\_CLASH\_INT} 中断的隐蔽中断状态位。（只读）
\item [UART\_RS485\_FRM\_ERR\_INT\_ST] \hyperref[int:UARTRS485FRMERRINT]{UART\_RS485\_FRM\_ERR\_INT} 中断的隐蔽中断状态位。（只读）
\item [UART\_RS485\_PARITY\_ERR\_INT\_ST] \hyperref[int:UARTRS485PARITYERRINT]{UART\_RS485\_PARITY\_ERR\_INT} 中断的隐蔽中断状态位。（只读）
\item [UART\_TX\_DONE\_INT\_ST] \hyperref[int:UARTTXDONEINT]{UART\_TX\_DONE\_INT} 中断的隐蔽中断状态位。（只读）
\item [UART\_TX\_BRK\_IDLE\_DONE\_INT\_ST] \hyperref[int:UARTTXBRKIDLEDONEINT]{UART\_TX\_BRK\_IDLE\_DONE\_INT} 中断的隐蔽中断状态位。（只读）
\item [UART\_TX\_BRK\_DONE\_INT\_ST] \hyperref[int:UARTTXBRKDONEINT]{UART\_TX\_BRK\_DONE\_INT} 中断的隐蔽中断状态位。（只读）
\item [UART\_GLITCH\_DET\_INT\_ST] \hyperref[int:UARTGLITCHDETINT]{UART\_GLITCH\_DET\_INT} 中断的隐蔽中断状态位。（只读）
\item [UART\_SW\_XOFF\_INT\_ST] \hyperref[int:UARTSWXOFFINT]{UART\_SW\_XOFF\_INT} 中断的隐蔽中断状态位。（只读）
\item [UART\_SW\_XON\_INT\_ST] \hyperref[int:UARTSWXONINT]{UART\_SW\_XON\_INT} 中断的隐蔽中断状态位。（只读）
\item [UART\_RXFIFO\_TOUT\_INT\_ST] \hyperref[int:UARTRXFIFOTOUTINT]{UART\_RXFIFO\_TOUT\_INT} 中断的隐蔽中断状态位。（只读）
\item [UART\_BRK\_DET\_INT\_ST] \hyperref[int:UARTBRKDETINT]{UART\_BRK\_DET\_INT} 中断的隐蔽中断状态位。（只读）
\item [UART\_CTS\_CHG\_INT\_ST] \hyperref[int:UARTCTSCHGINT]{UART\_CTS\_CHG\_INT} 中断的隐蔽中断状态位。（只读）
\item [UART\_DSR\_CHG\_INT\_ST] \hyperref[int:UARTDSRCHGINT]{UART\_DSR\_CHG\_INT} 中断的隐蔽中断状态位。（只读）
\item [UART\_RXFIFO\_OVF\_INT\_ST] \hyperref[int:UARTRXFIFOOVFINT]{UART\_RXFIFO\_OVF\_INT} 中断的隐蔽中断状态位。（只读）
\item [UART\_FRM\_ERR\_INT\_ST] \hyperref[int:UARTFRMERRINT]{UART\_FRM\_ERR\_INT} 中断的隐蔽中断状态位。（只读）
\item [UART\_PARITY\_ERR\_INT\_ST] \hyperref[int:UARTPARITYERRINT]{UART\_PARITY\_ERR\_INT} 中断的隐蔽中断状态位。（只读）
\item [UART\_TXFIFO\_EMPTY\_INT\_ST] \hyperref[int:UARTTXFIFOEMPTYINT]{UART\_TXFIFO\_EMPTY\_INT} 中断的隐蔽中断状态位。（只读）
\item [UART\_RXFIFO\_FULL\_INT\_ST] \hyperref[int:UARTRXFIFOFULLINT]{UART\_RXFIFO\_FULL\_INT} 中断的隐蔽中断状态位。（只读）
\end{reglist}\end{regdesc}
\end{register}


\begin{register}{H}{UART\_INT\_ENA\_REG}{0xC}\label{regdesc:UARTINTENAREG}
\regfield{(reserved)}{13}{19}{0000000000000}%
\regfield{UART\_AT\_CMD\_CHAR\_DET\_INT\_ENA}{1}{18}{0}%
\regfield{UART\_RS485\_CLASH\_INT\_ENA}{1}{17}{0}%
\regfield{UART\_RS485\_FRM\_ERR\_INT\_ENA}{1}{16}{0}%
\regfield{UART\_RS485\_PARITY\_ERR\_INT\_ENA}{1}{15}{0}%
\regfield{UART\_TX\_DONE\_INT\_ENA}{1}{14}{0}%
\regfield{UART\_TX\_BRK\_IDLE\_DONE\_INT\_ENA}{1}{13}{0}%
\regfield{UART\_TX\_BRK\_DONE\_INT\_ENA}{1}{12}{0}%
\regfield{UART\_GLITCH\_DET\_INT\_ENA}{1}{11}{0}%
\regfield{UART\_SW\_XOFF\_INT\_ENA}{1}{10}{0}%
\regfield{UART\_SW\_XON\_INT\_ENA}{1}{9}{0}%
\regfield{UART\_RXFIFO\_TOUT\_INT\_ENA}{1}{8}{0}%
\regfield{UART\_BRK\_DET\_INT\_ENA}{1}{7}{0}%
\regfield{UART\_CTS\_CHG\_INT\_ENA}{1}{6}{0}%
\regfield{UART\_DSR\_CHG\_INT\_ENA}{1}{5}{0}%
\regfield{UART\_RXFIFO\_OVF\_INT\_ENA}{1}{4}{0}%
\regfield{UART\_FRM\_ERR\_INT\_ENA}{1}{3}{0}%
\regfield{UART\_PARITY\_ERR\_INT\_ENA}{1}{2}{0}%
\regfield{UART\_TXFIFO\_EMPTY\_INT\_ENA}{1}{1}{0}%
\regfield{UART\_RXFIFO\_FULL\_INT\_ENA}{1}{0}{0}%
\reglabel{Reset}\regnewline%
\begin{regdesc}\begin{reglist}
\item [UART\_AT\_CMD\_CHAR\_DET\_INT\_ENA] \hyperref[int:UARTATCMDCHARDETINT]{UART\_AT\_CMD\_CHAR\_DET\_INT} 中断的使能位。（读/写）
\item [UART\_RS485\_CLASH\_INT\_ENA] \hyperref[int:UARTRS485CLASHINT]{UART\_RS485\_CLASH\_INT} 中断的使能位。（读/写）
\item [UART\_RS485\_FRM\_ERR\_INT\_ENA] \hyperref[int:UARTRS485FRMERRINT]{UART\_RS485\_FRM\_ERR\_INT} 中断的使能位。（读/写）
\item [UART\_RS485\_PARITY\_ERR\_INT\_ENA] \hyperref[int:UARTRS485PARITYERRINT]{UART\_RS485\_PARITY\_ERR\_INT} 中断的使能位。（读/写）
\item [UART\_TX\_DONE\_INT\_ENA]T \hyperref[int:UARTTXDONEINT]{UART\_TX\_DONE\_INT} 中断的使能位。（读/写）
\item [UART\_TX\_BRK\_IDLE\_DONE\_INT\_ENA] \hyperref[int:UARTTXBRKIDLEDONEINT]{UART\_TX\_BRK\_IDLE\_DONE\_INT} 中断的使能位。（读/写）
\item [UART\_TX\_BRK\_DONE\_INT\_ENA] \hyperref[int:UARTTXBRKDONEINT]{UART\_TX\_BRK\_DONE\_INT} 中断的使能位。（读/写）
\item [UART\_GLITCH\_DET\_INT\_ENA] \hyperref[int:UARTGLITCHDETINT]{UART\_GLITCH\_DET\_INT} 中断的使能位。（读/写）
\item [UART\_SW\_XOFF\_INT\_ENA] \hyperref[int:UARTSWXOFFINT]{UART\_SW\_XOFF\_INT} 中断的使能位。（读/写）
\item [UART\_SW\_XON\_INT\_ENA] \hyperref[int:UARTSWXONINT]{UART\_SW\_XON\_INT} 中断的使能位。（读/写）
\item [UART\_RXFIFO\_TOUT\_INT\_ENA] \hyperref[int:UARTRXFIFOTOUTINT]{UART\_RXFIFO\_TOUT\_INT} 中断的使能位。（读/写）
\item [UART\_BRK\_DET\_INT\_ENA] \hyperref[int:UARTBRKDETINT]{UART\_BRK\_DET\_INT} 中断的使能位。（读/写）
\item [UART\_CTS\_CHG\_INT\_ENA] \hyperref[int:UARTCTSCHGINT]{UART\_CTS\_CHG\_INT} 中断的使能位。（读/写）
\item [UART\_DSR\_CHG\_INT\_ENA] \hyperref[int:UARTDSRCHGINT]{UART\_DSR\_CHG\_INT} 中断的使能位。（读/写）
\item [UART\_RXFIFO\_OVF\_INT\_ENA] \hyperref[int:UARTRXFIFOOVFINT]{UART\_RXFIFO\_OVF\_INT} 中断的使能位。（读/写）
\item [UART\_FRM\_ERR\_INT\_ENA] \hyperref[int:UARTFRMERRINT]{UART\_FRM\_ERR\_INT} 中断的使能位。（读/写）
\item [UART\_PARITY\_ERR\_INT\_ENA] \hyperref[int:UARTPARITYERRINT]{UART\_PARITY\_ERR\_INT} 中断的使能位。（读/写）
\item [UART\_TXFIFO\_EMPTY\_INT\_ENA] \hyperref[int:UARTTXFIFOEMPTYINT]{UART\_TXFIFO\_EMPTY\_INT} 中断的使能位。（读/写）
\item [UART\_RXFIFO\_FULL\_INT\_ENA] \hyperref[int:UARTRXFIFOFULLINT]{UART\_RXFIFO\_FULL\_INT} 中断的使能位。（读/写）
\end{reglist}\end{regdesc}
\end{register}


\begin{register}{H}{UART\_INT\_CLR\_REG}{0x10}\label{regdesc:UARTINTCLRREG}
\regfield{(reserved)}{13}{19}{0000000000000}%
\regfield{UART\_AT\_CMD\_CHAR\_DET\_INT\_CLR}{1}{18}{0}%
\regfield{UART\_RS485\_CLASH\_INT\_CLR}{1}{17}{0}%
\regfield{UART\_RS485\_FRM\_ERR\_INT\_CLR}{1}{16}{0}%
\regfield{UART\_RS485\_PARITY\_ERR\_INT\_CLR}{1}{15}{0}%
\regfield{UART\_TX\_DONE\_INT\_CLR}{1}{14}{0}%
\regfield{UART\_TX\_BRK\_IDLE\_DONE\_INT\_CLR}{1}{13}{0}%
\regfield{UART\_TX\_BRK\_DONE\_INT\_CLR}{1}{12}{0}%
\regfield{UART\_GLITCH\_DET\_INT\_CLR}{1}{11}{0}%
\regfield{UART\_SW\_XOFF\_INT\_CLR}{1}{10}{0}%
\regfield{UART\_SW\_XON\_INT\_CLR}{1}{9}{0}%
\regfield{UART\_RXFIFO\_TOUT\_INT\_CLR}{1}{8}{0}%
\regfield{UART\_BRK\_DET\_INT\_CLR}{1}{7}{0}%
\regfield{UART\_CTS\_CHG\_INT\_CLR}{1}{6}{0}%
\regfield{UART\_DSR\_CHG\_INT\_CLR}{1}{5}{0}%
\regfield{UART\_RXFIFO\_OVF\_INT\_CLR}{1}{4}{0}%
\regfield{UART\_FRM\_ERR\_INT\_CLR}{1}{3}{0}%
\regfield{UART\_PARITY\_ERR\_INT\_CLR}{1}{2}{0}%
\regfield{UART\_TXFIFO\_EMPTY\_INT\_CLR}{1}{1}{0}%
\regfield{UART\_RXFIFO\_FULL\_INT\_CLR}{1}{0}{0}%
\reglabel{Reset}\regnewline%
\begin{regdesc}\begin{reglist}
\item [UART\_AT\_CMD\_CHAR\_DET\_INT\_CLR] 置位清除 \hyperref[int:UARTATCMDCHARDETINT]{UART\_AT\_CMD\_CHAR\_DET\_INT} 中断。（只写）
\item [UART\_RS485\_CLASH\_INT\_CLR] 置位清除 \hyperref[int:UARTRS485CLASHINT]{UART\_RS485\_CLASH\_INT} 中断。（只写）
\item [UART\_RS485\_FRM\_ERR\_INT\_CLR] 置位清除 \hyperref[int:UARTRS485FRMERRINT]{UART\_RS485\_FRM\_ERR\_INT} 中断。（只写）
\item [UART\_RS485\_PARITY\_ERR\_INT\_CLR] 置位清除 \hyperref[int:UARTRS485PARITYERRINT]{UART\_RS485\_PARITY\_ERR\_INT} 中断。（只写）
\item [UART\_TX\_DONE\_INT\_CLR] 置位清除 \hyperref[int:UARTTXDONEINT]{UART\_TX\_DONE\_INT} 中断。（只写）
\item [UART\_TX\_BRK\_IDLE\_DONE\_INT\_CLR] 置位清除 \hyperref[int:UARTTXBRKIDLEDONEINT]{UART\_TX\_BRK\_IDLE\_DONE\_INT} 中断。（只写）
\item [UART\_TX\_BRK\_DONE\_INT\_CLR] 置位清除 \hyperref[int:UARTTXBRKDONEINT]{UART\_TX\_BRK\_DONE\_INT} 中断。（只写）
\item [UART\_GLITCH\_DET\_INT\_CLR] 置位清除 \hyperref[int:UARTGLITCHDETINT]{UART\_GLITCH\_DET\_INT} 中断。（只写）
\item [UART\_SW\_XOFF\_INT\_CLR] 置位清除 \hyperref[int:UARTSWXOFFINT]{UART\_SW\_XOFF\_INT} 中断。（只写）
\item [UART\_SW\_XON\_INT\_CLR] 置位清除 \hyperref[int:UARTSWXONINT]{UART\_SW\_XON\_INT} 中断。（只写）
\item [UART\_RXFIFO\_TOUT\_INT\_CLR]\label{UARTRXFIFOTOUTINTCLR} 置位清除 \hyperref[int:UARTRXFIFOTOUTINT]{UART\_RXFIFO\_TOUT\_INT} 中断。该寄存器只能在 rxfifo\_cnt 和 rx\_mem\_cnt 都为 0 的情况下才可以置位。（只写）
\item [UART\_BRK\_DET\_INT\_CLR] 置位清除 \hyperref[int:UARTBRKDETINT]{UART\_BRK\_DET\_INT} 中断。（只写）
\item [UART\_CTS\_CHG\_INT\_CLR] 置位清除 \hyperref[int:UARTCTSCHGINT]{UART\_CTS\_CHG\_INT} 中断。（只写）
\item [UART\_DSR\_CHG\_INT\_CLR] 置位清除 \hyperref[int:UARTDSRCHGINT]{UART\_DSR\_CHG\_INT} 中断。（只写）
\item [UART\_RXFIFO\_OVF\_INT\_CLR] 置位清除 \hyperref[int:UARTRXFIFOOVFINT]{UART\_RXFIFO\_OVF\_INT} 中断。（只写）
\item [UART\_FRM\_ERR\_INT\_CLR] 置位清除 \hyperref[int:UARTFRMERRINT]{UART\_FRM\_ERR\_INT} 中断。（只写）
\item [UART\_PARITY\_ERR\_INT\_CLR] 置位清除 \hyperref[int:UARTPARITYERRINT]{UART\_PARITY\_ERR\_INT} 中断。（只写）
\item [UART\_TXFIFO\_EMPTY\_INT\_CLR] 置位清除 \hyperref[int:UARTTXFIFOEMPTYINT]{UART\_TXFIFO\_EMPTY\_INT} 中断。（只写）
\item [UART\_RXFIFO\_FULL\_INT\_CLR]\label{UARTRXFIFOFULLINTCLR} 置位清除 \hyperref[int:UARTRXFIFOFULLINT]{UART\_RXFIFO\_FULL\_INT} 中断。该寄存器只有在 Rx\_FIFO 中数据个数小于 \hyperref[UARTRXFIFOFULLTHRHD]{UART\_RXFIFO\_FULL\_THRHD} 的情况下置位才能生效。（只写）
\end{reglist}\end{regdesc}
\end{register}


\begin{register}{H}{UART\_CLKDIV\_REG}{0x14}\label{regdesc:UARTCLKDIVREG}
\regfield{(reserved)}{8}{24}{00000000}%
\regfield{UART\_CLKDIV\_FRAG}{4}{20}{{0x00}}%
\regfield{UART\_CLKDIV}{20}{0}{{0x0002B6}}%
\reglabel{Reset}\regnewline%
\begin{regdesc}\begin{reglist}
\item [UART\_CLKDIV\_FRAG] 分频系数的小数部分。（读/写）
\item [UART\_CLKDIV] 分频系数的整数部分。（读/写）
\end{reglist}\end{regdesc}
\end{register}


\begin{register}{H}{UART\_AUTOBAUD\_REG}{0x18}\label{regdesc:UARTAUTOBAUDREG}
\regfield{(reserved)}{16}{16}{0000000000000000}%
\regfield{UART\_GLITCH\_FILT}{8}{8}{{0x010}}%
\regfield{(reserved)}{7}{1}{0000000}%
\regfield{UART\_AUTOBAUD\_EN}{1}{0}{0}%
\reglabel{Reset}\regnewline%
\begin{regdesc}\begin{reglist}
\item [UART\_GLITCH\_FILT] 滤波门限值，当输入脉冲宽度小于此寄存器的值时，脉冲被忽略。此寄存器用于自动波特率检测的过程中。（读/写）
\item [UART\_AUTOBAUD\_EN] 自动波特率检测的使能位。（读/写）
\end{reglist}\end{regdesc}
\end{register}


\begin{register}{H}{UART\_STATUS\_REG}{0x1C}\label{regdesc:UARTSTATUSREG}
\regfield{UART\_TXD}{1}{31}{{0x000}}%
\regfield{UART\_RTSN}{1}{30}{0}%
\regfield{UART\_DTRN}{1}{29}{0}%
\regfield{(reserved)}{1}{28}{0}%
\regfield{UART\_ST\_UTX\_OUT}{4}{24}{0000}%
\regfield{UART\_TXFIFO\_CNT}{8}{16}{00000000}%
\regfield{UART\_RXD}{1}{15}{0}%
\regfield{UART\_CTSN}{1}{14}{0}%
\regfield{UART\_DSRN}{1}{13}{0}%
\regfield{(reserved)}{1}{12}{0}%
\regfield{UART\_ST\_URX\_OUT}{4}{8}{0000}%
\regfield{UART\_RXFIFO\_CNT}{8}{0}{00000000}%
\reglabel{Reset}\regnewline%
\begin{regdesc}\begin{reglist}
\item [UART\_TXD] 此位表明内部 UART RxD 信号的电平。（只读）
\item [UART\_RTSN] 此位对应内部 UART CTS 信号的电平。（只读）
\item [UART\_DTRN]  此位对应内部 UAR DSR 信号的电平。（只读）
\item [UART\_ST\_UTX\_OUT] 此寄存器存储发送器有限状态机的状态。0: TX\_IDLE；1: TX\_STRT；2: TX\_DAT0；3: TX\_DAT1；4: TX\_DAT2；5: TX\_DAT3；6: TX\_DAT4；7: TX\_DAT5；8: TX\_DAT6；9: TX\_DAT7；10: TX\_PRTY；11: TX\_STP1；12: TX\_STP2；13: TX\_DL0；14: TX\_DL1。（只读）
\item [UART\_TXFIFO\_CNT]\label{UARTTXFIFOCNT} (tx\_mem\_cnt, txfifo\_cnt) 存储发送 FIFO 中的有效数据字节数。 tx\_mem\_cnt 存储 3 个最高位；txfifo\_cnt 存储 8 个最低位。（只读）
\item [UART\_RXD] 此位对应内部 UART RxD 信号的电平。（只读）
\item [UART\_CTSN] 此位对应内部 UART CTS 信号的电平。（只读）
\item [UART\_DSRN] 此位对应内部 UAR DSR 信号的电平。（只读）
\item [UART\_ST\_URX\_OUT] 此寄存器存储接收器有限状态机的状态。 0: RX\_IDLE；1: RX\_STRT；2: RX\_DAT0；3: RX\_DAT1； 4: RX\_DAT2；5: RX\_DAT3；6: RX\_DAT4；7: RX\_DAT5；8: RX\_DAT6；9: RX\_DAT7；10: RX\_PRTY；11: RX\_STP1；12: RX\_STP2；13: RX\_DL1。（只读）
\item [UART\_RXFIFO\_CNT]\label{UARTRXFIFOCNT} (rx\_mem\_cnt, rxfifo\_cnt) 存储接收 FIFO 中的有效数据字节数。 rx\_mem\_cnt 存储 3 个最高位；rxfifo\_cnt 存储 8 个最低位。（只读）
\end{reglist}\end{regdesc}
\end{register}


\begin{register}{H}{UART\_CONF0\_REG}{0x20}\label{regdesc:UARTCONF0REG}
\regfield{(reserved)}{4}{28}{0000}%
\regfield{UART\_TICK\_REF\_ALWAYS\_ON}{1}{27}{1}%
\regfield{(reserved)}{2}{25}{00}%
\regfield{UART\_DTR\_INV}{1}{24}{0}%
\regfield{UART\_RTS\_INV}{1}{23}{0}%
\regfield{UART\_TXD\_INV}{1}{22}{0}%
\regfield{UART\_DSR\_INV}{1}{21}{0}%
\regfield{UART\_CTS\_INV}{1}{20}{0}%
\regfield{UART\_RXD\_INV}{1}{19}{0}%
\regfield{UART\_TXFIFO\_RST}{1}{18}{0}%
\regfield{UART\_RXFIFO\_RST}{1}{17}{0}%
\regfield{UART\_IRDA\_EN}{1}{16}{0}%
\regfield{UART\_TX\_FLOW\_EN}{1}{15}{0}%
\regfield{UART\_LOOPBACK}{1}{14}{0}%
\regfield{UART\_IRDA\_RX\_INV}{1}{13}{0}%
\regfield{UART\_IRDA\_TX\_INV}{1}{12}{0}%
\regfield{UART\_IRDA\_WCTL}{1}{11}{0}%
\regfield{UART\_IRDA\_TX\_EN}{1}{10}{0}%
\regfield{UART\_IRDA\_DPLX}{1}{9}{0}%
\regfield{UART\_TXD\_BRK}{1}{8}{0}%
\regfield{UART\_SW\_DTR}{1}{7}{0}%
\regfield{UART\_SW\_RTS}{1}{6}{0}%
\regfield{UART\_STOP\_BIT\_NUM}{2}{4}{{1}}%
\regfield{UART\_BIT\_NUM}{2}{2}{{3}}%
\regfield{UART\_PARITY\_EN}{1}{1}{0}%
\regfield{UART\_PARITY}{1}{0}{0}%
\reglabel{Reset}\regnewline%
\begin{regdesc}\begin{reglist}
\item [UART\_TICK\_REF\_ALWAYS\_ON] 此寄存器用于选择时钟。1：APB 时钟；0: REF\_TICK。（读/写）
\item [UART\_DTR\_INV] 置位反转 UART DTR 信号的电平。（读/写）
\item [UART\_RTS\_INV] 置位反转 UART RTS 信号的电平。（读/写）
\item [UART\_TXD\_INV] 置位反转 UART TxD 信号的电平。（读/写）
\item [UART\_DSR\_INV] 置位反转 UART DSR 信号的电平。（读/写）
\item [UART\_CTS\_INV] 置位反转 UART CTS 信号的电平。（读/写）
\item [UART\_RXD\_INV] 置位反转 UART Rxd 信号的电平。（读/写）
\item [UART\_TXFIFO\_RST] 置位，复位 UART 发送 FIFO。注意，UART2 没有该复位寄存器，且 UART1 的 UART1\_TXFIFO\_RST 和 UART1\_RXFIFO\_RST 会影响 UART2 的工作。因此，只有在 UART2 的 Tx\_FIFO 和 Rx\_FIFO 中没有数据时，才可以置位这两个寄存器 。（读/写）
\item [UART\_RXFIFO\_RST] 置位，复位 UART 接收 FIFO。注意，UART2 没有该复位寄存器，且 UART1 的 UART1\_TXFIFO\_RST 和 UART1\_RXFIFO\_RST 会影响 UART2 的工作。因此，只有在 UART2 的 Tx\_FIFO 和 Rx\_FIFO 中没有数据时，才可以置位这两个寄存器 。（读/写）
\item [UART\_IRDA\_EN] 置位使能 IrDA 协议。（读/写）
\item [UART\_TX\_FLOW\_EN] 置位使能发送器的流控功能。（读/写）
\item [UART\_LOOPBACK] 置位使能 UART 回环测试功能。（读/写）
\item [UART\_IRDA\_RX\_INV] 置位反转 IrDA 接收器的电平。（读/写）
\item [UART\_IRDA\_TX\_INV] 置位反转 IrDA 发送器的电平。（读/写）
\item [UART\_IRDA\_WCTL]1：IrDA 发送器的第 11 位与第 10 位相同。0:设置 IrDA 发送器的第 11 位为 0。（读/写）
\item [UART\_IRDA\_TX\_EN] IrDA 发送器的启动使能位。（读/写）
\item [UART\_IRDA\_DPLX] 置位使能 IrDA 回环模式。（读/写）
\item [UART\_TXD\_BRK] 置位使发送器在数据发送完成后发送 NULL。（读/写）
\item [UART\_SW\_DTR] 配置用于软件流控的软件 DTR 信号。（读/写）
\item [UART\_SW\_RTS] 用于 UART\_RX\_FLOW\_EN 为 0 时的硬件流控。置位拉低 RTS (rtsn\_out) 信号，复位拉高信号。（读/写）
\item [寄存器描述下一页继续。]
\end{reglist}\end{regdesc}
\end{register}
\addtocounter{Regfloat}{-1}
\begin{register}{H}{UART\_CONF0\_REG}{0x20}
\begin{regdesc}\begin{reglist}
\item [继上一页寄存器描述。]
\item [UART\_STOP\_BIT\_NUM] 用于设置停止位的长度。 \\0: 无效值，没有作用\\
1: 1 位\\
2: 1.5 位\\
3: 2 位\\（读/写）
\item [UART\_BIT\_NUM] 用于设置数据的长度；0: 5 bit；1: 6 bit；2: 7 bit；3: 8 bit。（读/写）
\item [UART\_PARITY\_EN] 置位使能 UART 奇偶校验。（读/写）
\item [UART\_PARITY] 配置奇偶校验方式。0: 偶校验；1: 奇校验。（读/写）
\end{reglist}\end{regdesc}
\end{register}

\begin{register}{H}{UART\_CONF1\_REG}{0x24}\label{regdesc:UARTCONF1REG}
\regfield{UART\_RX\_TOUT\_EN}{1}{31}{0}%
\regfield{UART\_RX\_TOUT\_THRHD}{7}{24}{0000000}%
\regfield{UART\_RX\_FLOW\_EN}{1}{23}{0}%
\regfield{UART\_RX\_FLOW\_THRHD}{7}{16}{{0x00}}%
\regfield{(reserved)}{1}{15}{0}%
\regfield{UART\_TXFIFO\_EMPTY\_THRHD}{7}{8}{{0x60}}%
\regfield{(reserved)}{1}{7}{0}%
\regfield{UART\_RXFIFO\_FULL\_THRHD}{7}{0}{{0x60}}%
\reglabel{Reset}\regnewline%
\begin{regdesc}\begin{reglist}
\item [UART\_RX\_TOUT\_EN] 置位使能 UART 接收器的超时功能。（读/写）
\item [UART\_RX\_TOUT\_THRHD]\label{UARTRXTOUTTHRHD} 配置 UART 接收器等待接收的超时时间。在使用 APB\_CLK 时钟源时，该寄存器以波特率周期的 8 倍 为计时单位；在使用 REF\_TICK 时钟源时，该寄存器的计时单位为 UART 波特率周期 * 8 * (REF\_TICK 频率)/(APB\_CLK 频率)。（读/写）
\item [UART\_RX\_FLOW\_EN] 置位使能 UART 接收器的流控功能。1: 配置 sw\_rts 信号选择软件流控；0：关闭软件流控。（读/写）
\item [UART\_RX\_FLOW\_THRHD] 当 UART\_RX\_FLOW\_EN 为 1、接收 FIFO 超过阈值时，接收器产生 rtsn\_out 信号告诉发送器停止发送数据。阈值为 (rx\_flow\_thrhd\_h3, rx\_flow\_thrhd)。（读/写）
\item [UART\_TXFIFO\_EMPTY\_THRHD]\label{UARTTXFIFOEMPTYTHRHD} 当发送 FIFO 的数据量少于阈值时，会产生 TXFIFO\_EMPTY\_INT\_RAW 中断。阈值为 (tx\_mem\_empty\_thrhd, txfifo\_empty\_thrhd)。（读/写）
\item [UART\_RXFIFO\_FULL\_THRHD]\label{UARTRXFIFOFULLTHRHD} 当接收器接收到比阈值多的数据时，接收器产生 RXFIFO\_FULL\_INT\_RAW 中断。阈值为 (rx\_flow\_thrhd\_h3, rxfifo\_full\_thrhd)。（读/写）
\end{reglist}\end{regdesc}
\end{register}


\begin{register}{H}{UART\_LOWPULSE\_REG}{0x28}\label{regdesc:UARTLOWPULSEREG}
\regfield{(reserved)}{12}{20}{000000000000}%
\regfield{UART\_LOWPULSE\_MIN\_CNT}{20}{0}{{0x0FFFFF}}%
\reglabel{Reset}\regnewline%
\begin{regdesc}\begin{reglist}
\item [UART\_LOWPULSE\_MIN\_CNT] 此寄存器存储最小低电平脉冲宽度，用于波特率自检过程。（只读）
\end{reglist}\end{regdesc}
\end{register}


\begin{register}{H}{UART\_HIGHPULSE\_REG}{0x2C}\label{regdesc:UARTHIGHPULSEREG}
\regfield{(reserved)}{12}{20}{000000000000}%
\regfield{UART\_HIGHPULSE\_MIN\_CNT}{20}{0}{{0x0FFFFF}}%
\reglabel{Reset}\regnewline%
\begin{regdesc}\begin{reglist}
\item [UART\_HIGHPULSE\_MIN\_CNT] 此寄存器存储最小高电平脉冲宽度值，用于波特率自检过程。（只读）
\end{reglist}\end{regdesc}
\end{register}


\begin{register}{H}{UART\_RXD\_CNT\_REG}{0x30}\label{regdesc:UARTRXDCNTREG}
\regfield{(reserved)}{22}{10}{0000000000000000000000}%
\regfield{UART\_RXD\_EDGE\_CNT}{10}{0}{{0x000}}%
\reglabel{Reset}\regnewline%
\begin{regdesc}\begin{reglist}
\item [UART\_RXD\_EDGE\_CNT] 此寄存器存储 RxD 沿变化的次数，用于波特率自检过程。（只读）
\end{reglist}\end{regdesc}
\end{register}


\begin{register}{H}{UART\_FLOW\_CONF\_REG}{0x34}\label{regdesc:UARTFLOWCONFREG}
\regfield{(reserved)}{26}{6}{00000000000000000000000000}%
\regfield{UART\_SEND\_XOFF}{1}{5}{0}%
\regfield{UART\_SEND\_XON}{1}{4}{0}%
\regfield{UART\_FORCE\_XOFF}{1}{3}{0}%
\regfield{UART\_FORCE\_XON}{1}{2}{0}%
\regfield{UART\_XONOFF\_DEL}{1}{1}{0}%
\regfield{UART\_SW\_FLOW\_CON\_EN}{1}{0}{0}%
\reglabel{Reset}\regnewline%
\begin{regdesc}\begin{reglist}
\item [UART\_SEND\_XOFF]\label{regdesc:UARTSENDXOFF} 硬件自清 0，置位发送 Xoff 字符。（读/写）
\item [UART\_SEND\_XON]\label{regdesc:UARTSENDXON} 硬件自清 0，置位发送 Xon 字符。（读/写）
\item [UART\_FORCE\_XOFF]\label{regdesc:UARTFORCEXOFF} 置位设置 CTSn 阻止发送器发送数据。（读/写）
\item [UART\_FORCE\_XON]\label{regdesc:UARTFORCEXON} 置位清除 CTSn 使能发送器继续发送数据。（读/写）
\item [UART\_XONOFF\_DEL] 置位移除接收到的数据中的流控字符。（读/写）
\item [UART\_SW\_FLOW\_CON\_EN] 置位使能软件流控，与 sw\_xon 或 sw\_xoff 寄存器一起使用。（读/写）
\end{reglist}\end{regdesc}
\end{register}


\begin{register}{H}{UART\_SLEEP\_CONF\_REG}{0x38}\label{regdesc:UARTSLEEPCONFREG}
\regfield{(reserved)}{22}{10}{0000000000000000000000}%
\regfield{UART\_ACTIVE\_THRESHOLD}{10}{0}{{0x0F0}}%
\reglabel{Reset}\regnewline%
\begin{regdesc}\begin{reglist}
\item [UART\_ACTIVE\_THRESHOLD]\label{UARTACTIVETHRESHOLD} 当输入 RxD 沿变化的次数大于等于 (UART\_ACTIVE\_THRESHOLD+2) 时，系统从 Light-sleep 中醒来。（读/写）
\end{reglist}\end{regdesc}
\end{register}


\begin{register}{H}{UART\_SWFC\_CONF\_REG}{0x3C}\label{regdesc:UARTSWFCCONFREG}
\regfield{UART\_XOFF\_CHAR}{8}{24}{{0x013}}%
\regfield{UART\_XON\_CHAR}{8}{16}{{0x011}}%
\regfield{UART\_XOFF\_THRESHOLD}{8}{8}{{0x0E0}}%
\regfield{UART\_XON\_THRESHOLD}{8}{0}{{0x000}}%
\reglabel{Reset}\regnewline%
\begin{regdesc}\begin{reglist}
\item [UART\_XOFF\_CHAR] 存储 Xoff 流控字符。（读/写）
\item [UART\_XON\_CHAR] 存储 Xon 流控字符。（读/写）
\item [UART\_XOFF\_THRESHOLD]\label{UARTXOFFTHRESHOLD} 当接收 FIFO 中的数据量多于配置值时，会发送一个 Xoff 字符，需要 UART\_SW\_FLOW\_CON\_EN 置为 1。（读/写）
\item [UART\_XON\_THRESHOLD]\label{UARTXONTHRESHOLD} 当接收 FIFO 中的数据量少于配置值时，会发送一个 Xon 字符，需要 UART\_SW\_FLOW\_CON\_EN 置为 1。（读/写）
\end{reglist}\end{regdesc}
\end{register}


\begin{register}{H}{UART\_IDLE\_CONF\_REG}{0x40}\label{regdesc:UARTIDLECONFREG}
\regfield{(reserved)}{4}{28}{0000}%
\regfield{UART\_TX\_BRK\_NUM}{8}{20}{{0x00A}}%
\regfield{UART\_TX\_IDLE\_NUM}{10}{10}{{0x100}}%
\regfield{UART\_RX\_IDLE\_THRHD}{10}{0}{{0x100}}%
\reglabel{Reset}\regnewline%
\begin{regdesc}\begin{reglist}
\item [UART\_TX\_BRK\_NUM] 用于在发送数据结束后配置发送的 0 的数量，当 txd\_brk 置为 1 时工作。（读/写）
\item [UART\_TX\_IDLE\_NUM] 用于配置数据传输的间隔。（读/写）
\item [UART\_RX\_IDLE\_THRHD] 当接收器等待的时间大于配置值时，会产生帧结束信号。（读/写）
\end{reglist}\end{regdesc}
\end{register}


\begin{register}{H}{UART\_RS485\_CONF\_REG}{0x44}\label{regdesc:UARTRS485CONFREG}
\regfield{(reserved)}{22}{10}{0000000000000000000000}%
\regfield{UART\_RS485\_TX\_DLY\_NUM}{4}{6}{0000}%
\regfield{UART\_RS485\_RX\_DLY\_NUM}{1}{5}{0}%
\regfield{UART\_RS485RXBY\_TX\_EN}{1}{4}{0}%
\regfield{UART\_RS485TX\_RX\_EN}{1}{3}{0}%
\regfield{UART\_DL1\_EN}{1}{2}{0}%
\regfield{UART\_DL0\_EN}{1}{1}{0}%
\regfield{UART\_RS485\_EN}{1}{0}{0}%
\reglabel{Reset}\regnewline%
\begin{regdesc}\begin{reglist}
\item [UART\_RS485\_TX\_DLY\_NUM] 用于延迟发送器内部数据信号。（读/写）
\item [UART\_RS485\_RX\_DLY\_NUM] 用于延迟接收器内部数据信号。（读/写）
\item [UART\_RS485RXBY\_TX\_EN]1: 当 RS-485 接收器线路繁忙时 RS-485 发送器能够发送数据。0: 当接收器繁忙时 RS-485 发送器不发送数据。（读/写）
\item [UART\_RS485TX\_RX\_EN] 置位使能发送器输出信号回环到接收器输入信号。（读/写）
\item [UART\_DL1\_EN] 置位延迟停止位 1 bit。（读/写）
\item [UART\_DL0\_EN]\label{UARTDL0EN} 置位之后，在 DL1 之后延迟 1 bit 停止位。（读/写）
\item [UART\_RS485\_EN] 置位选择 RS-485 模式。（读/写）
\end{reglist}\end{regdesc}
\end{register}


\begin{register}{H}{UART\_AT\_CMD\_PRECNT\_REG}{0x48}\label{regdesc:UARTATCMDPRECNTREG}
\regfield{(reserved)}{8}{24}{00000000}%
\regfield{UART\_PRE\_IDLE\_NUM}{24}{0}{{0x0186A00}}%
\reglabel{Reset}\regnewline%
\begin{regdesc}\begin{reglist}
\item [UART\_PRE\_IDLE\_NUM] 用于配置接收器接收到第一个 at\_cmd 前的空闲时间。当空闲时间少于配置值时，接收器不会将接收到的下一个数据当作 at\_cmd 字符。（读/写）
\end{reglist}\end{regdesc}
\end{register}


\begin{register}{H}{UART\_AT\_CMD\_POSTCNT\_REG}{0x4c}\label{regdesc:UARTATCMDPOSTCNTREG}
\regfield{(reserved)}{8}{24}{00000000}%
\regfield{UART\_POST\_IDLE\_NUM}{24}{0}{{0x0186A00}}%
\reglabel{Reset}\regnewline%
\begin{regdesc}\begin{reglist}
\item [UART\_POST\_IDLE\_NUM] 用于配置最后一个 at\_cmd 和下一个数据之间的间隔时间。当间隔时间少于配置值时，不会将前一个数据当做 at\_cmd 字符。（读/写）
\end{reglist}\end{regdesc}
\end{register}


\begin{register}{H}{UART\_AT\_CMD\_GAPTOUT\_REG}{0x50}\label{regdesc:UARTATCMDGAPTOUTREG}
\regfield{(reserved)}{8}{24}{00000000}%
\regfield{UART\_RX\_GAP\_TOUT}{24}{0}{{0x0001E00}}%
\reglabel{Reset}\regnewline%
\begin{regdesc}\begin{reglist}
\item [UART\_RX\_GAP\_TOUT]\label{UARTRXGAPTOUT} 用于配置 at\_cmd 字符之间的间隔时间。当间隔时间大于配置值时，不会将数据当做连续的 at\_cmd 字符，且该寄存器的配置值至少要大于二分之一波特率。（读/写）
\end{reglist}\end{regdesc}
\end{register}


\begin{register}{H}{UART\_AT\_CMD\_CHAR\_REG}{0x54}\label{regdesc:UARTATCMDCHARREG}
\regfield{(reserved)}{16}{16}{0000000000000000}%
\regfield{UART\_CHAR\_NUM}{8}{8}{{0x003}}%
\regfield{UART\_AT\_CMD\_CHAR}{8}{0}{{0x02B}}%
\reglabel{Reset}\regnewline%
\begin{regdesc}\begin{reglist}
\item [UART\_CHAR\_NUM] 用于配置接收器接收到的连续 at\_cmd 字符的数量。（读/写）
\item [UART\_AT\_CMD\_CHAR] 用于配置 at\_cmd 字符的内容。（读/写）
\end{reglist}\end{regdesc}
\end{register}


\begin{register}{H}{UART\_MEM\_CONF\_REG}{0x58}\label{regdesc:UARTMEMCONFREG}
\regfield{(reserved)}{1}{31}{0}%
\regfield{UART\_TX\_MEM\_EMPTY\_THRHD}{3}{28}{{0x0}}%
\regfield{UART\_RX\_MEM\_FULL\_THRHD}{3}{25}{{0x0}}%
\regfield{UART\_XOFF\_THRESHOLD\_H2}{2}{23}{{0x0}}%
\regfield{UART\_XON\_THRESHOLD\_H2}{2}{21}{{0x0}}%
\regfield{UART\_RX\_TOUT\_THRHD\_H3}{3}{18}{{0x0}}%
\regfield{UART\_RX\_FLOW\_THRHD\_H3}{3}{15}{{0x0}}%
\regfield{(reserved)}{4}{11}{0000}%
\regfield{UART\_TX\_SIZE}{4}{7}{{0x01}}%
\regfield{UART\_RX\_SIZE}{4}{3}{{0x01}}%
\regfield{(reserved)}{2}{1}{00}%
\regfield{UART\_MEM\_PD}{1}{0}{0}%
\reglabel{Reset}\regnewline%
\begin{regdesc}\begin{reglist}
\item [UART\_TX\_MEM\_EMPTY\_THRHD] 参考 \hyperref[UARTTXFIFOEMPTYTHRHD]{TXFIFO\_EMPTY\_THRHD} 的描述。（读/写）
\item [UART\_RX\_MEM\_FULL\_THRHD] 参考 \hyperref[UARTRXFIFOFULLTHRHD]{RXFIFO\_FULL\_THRHD} 的描述。（读/写）
\item [UART\_XOFF\_THRESHOLD\_H2] 参考 \hyperref[UARTXOFFTHRESHOLD]{UART\_XOFF\_THRESHOLD} 的描述。（读/写）
\item [UART\_XON\_THRESHOLD\_H2] 参考 \hyperref[UARTXONTHRESHOLD]{UART\_XON\_THRESHOLD} 的描述。（读/写）
\item [UART\_RX\_TOUT\_THRHD\_H3] 参考 \hyperref[UARTRXTOUTTHRHD]{RX\_TOUT\_THRHD} 的描述。（读/写）
\item [UART\_RX\_FLOW\_THRHD\_H3] 参考 \hyperref[UARTRXFIFOFULLTHRHD]{RX\_FLOW\_THRHD} 的描述。（读/写）
\item [UART\_TX\_SIZE] 用于配置分配给发送 FIFO 的 RAM 空间，默认为 128 byte。（读/写）
\item [UART\_RX\_SIZE]用于配置分配给接收 FIFO 的 RAM 空间，默认为 128 byte。（读/写）
\item [UART\_MEM\_PD] 置位关闭 RAM，当 3 个 UART 控制器的 reg\_mem\_pd 都置为 1 时，RAM 进入低功耗模式。（读/写）
\end{reglist}\end{regdesc}
\end{register}


\begin{register}{H}{UART\_MEM\_TX\_STATUS\_REG}{0x5c}\label{regdesc:UARTMEMTXSTATUSREG}
\regfield{(reserved)}{8}{24}{0}%
\regfield{UART\_MEM\_TX\_WR\_ADDR}{11}{13}{{0}}%
\regfield{UART\_MEM\_TX\_RD\_ADDR}{11}{2}{0}%
\regfield{(reserved)}{2}{0}{{0}}%
\reglabel{Reset}\regnewline%
\begin{regdesc}\begin{reglist}
\item [UART\_MEM\_TX\_WR\_ADDR] 表示写 TX FIFO 的偏移地址。 (RO)
\item [UART\_MEM\_TX\_RD\_ADDR] 表示读 TX FIFO 的偏移地址。 (RO)
\end{reglist}\end{regdesc}
\end{register}


\begin{register}{H}{UART\_MEM\_RX\_STATUS\_REG}{0x0060}\label{regdesc:UARTMEMRXSTATUSREG}
\regfield{(reserved)}{8}{24}{0}%
\regfield{UART\_MEM\_RX\_WR\_ADDR}{11}{13}{{0}}%
\regfield{UART\_MEM\_RX\_RD\_ADDR}{11}{2}{0}%
\regfield{(reserved)}{2}{0}{{0}}%
\reglabel{Reset}\regnewline%
\begin{regdesc}\begin{reglist}
\item [UART\_MEM\_RX\_RD\_ADDR] 表示读 RX FIFO 的偏移地址。 (RO)
\item [UART\_MEM\_RX\_WR\_ADDR] 表示写 RX FIFO 的偏移地址。 (RO)
\end{reglist}\end{regdesc}
\end{register}


\begin{register}{H}{UART\_MEM\_CNT\_STATUS\_REG}{0x64}\label{regdesc:UARTMEMCNTSTATUSREG}
\regfield{(reserved)}{26}{6}{00000000000000000000000000}%
\regfield{UART\_TX\_MEM\_CNT}{3}{3}{000}%
\regfield{UART\_RX\_MEM\_CNT}{3}{0}{000}%
\reglabel{Reset}\regnewline%
\begin{regdesc}\begin{reglist}
\item [UART\_TX\_MEM\_CNT] 参考 \hyperref[UARTTXFIFOCNT]{TXFIFO\_CNT} 的描述。（只读）
\item [UART\_RX\_MEM\_CNT] 参考 \hyperref[UARTRXFIFOCNT]{RXFIFO\_CNT} 的描述。（只读）
\end{reglist}\end{regdesc}
\end{register}


\begin{register}{H}{UART\_POSPULSE\_REG}{0x68}\label{regdesc:UARTPOSPULSEREG}
\regfield{(reserved)}{12}{20}{000000000000}%
\regfield{UART\_POSEDGE\_MIN\_CNT}{20}{0}{{0x0FFFFF}}%
\reglabel{Reset}\regnewline%
\begin{regdesc}\begin{reglist}
\item [UART\_POSEDGE\_MIN\_CNT] 存储 RxD 上升沿的沿变化次数，用于波特率自检的过程。（只读）
\end{reglist}\end{regdesc}
\end{register}


\begin{register}{H}{UART\_NEGPULSE\_REG}{0x6c}\label{regdesc:UARTNEGPULSEREG}
\regfield{(reserved)}{12}{20}{000000000000}%
\regfield{UART\_NEGEDGE\_MIN\_CNT}{20}{0}{{0x0FFFFF}}%
\reglabel{Reset}\regnewline%
\begin{regdesc}\begin{reglist}
\item [UART\_NEGEDGE\_MIN\_CNT] 存储 RxD 下降沿的沿变化次数，用于波特率自检的过程。（只读）
\end{reglist}\end{regdesc}
\end{register}
